\chapter{Geriatrician}
\index{Geriatrician}

\section*{Definition and Overview}
A geriatrician is a medical specialist who cares for older people, particularly those with complex health needs. Geriatricians have completed medical school, specialist physician training, and additional advanced training in the medicine of ageing. They focus on the prevention, diagnosis, and management of illness in older adults, taking into account the interaction between physical, mental, social, and functional factors. Geriatricians lead multidisciplinary teams that include nurses, physiotherapists, occupational therapists, pharmacists, and social workers, working together to help older people maintain independence and quality of life. They are experts in falls, frailty, dementia, incontinence, polypharmacy, and the care of people with multiple coexisting conditions.

\section*{Epidemiology}
The need for geriatricians is growing rapidly as populations age. In many countries, the proportion of people aged 65 and over is increasing steadily, and those aged 85 and over are the fastest-growing age group. Older people account for a large share of hospital admissions, and many have multiple chronic conditions requiring coordinated care. There is a well-documented shortage of geriatricians relative to demand, both in many countries and internationally. Geriatric medicine is therefore focused on the most complex patients, while GPs, nurses, and other specialists provide most day-to-day care of older people in the community.

\section*{Aetiology and Pathophysiology}
This chapter concerns the role of a medical specialist rather than a disease process. The core skills of a geriatrician include:

\begin{itemize}
    \item \textbf{Comprehensive geriatric assessment:} a multidimensional evaluation of medical, functional, psychological, and social issues
    \item \textbf{Management of multimorbidity:} balancing the treatment of several chronic conditions and their medicines
    \item \textbf{Frailty identification:} recognising the syndrome of reduced physiological reserve that predicts poor outcomes
    \item \textbf{Cognition and mood:} assessment and management of dementia, delirium, and depression
    \item \textbf{Mobility and falls:} investigating the causes of falls and preventing recurrence
    \item \textbf{Continence:} assessing and treating urinary and faecal incontinence
    \item \textbf{Polypharmacy:} reviewing and simplifying medicines to reduce harm
    \item \textbf{End-of-life care:} advance care planning and palliative approaches in advanced age
\end{itemize}

\section*{Clinical Features}
People are usually referred to a geriatrician when they have features suggesting that their care needs a specialist view.

\begin{itemize}
    \item Multiple chronic conditions that are difficult to manage together
    \item Recurrent falls or a decline in mobility and balance
    \item Memory problems, confusion, or suspected dementia
    \item Sudden or gradual functional decline, such as difficulty with daily tasks
    \item Complex medicine regimens with side effects or interactions
    \item Frailty, weight loss, and loss of strength
    \item Incontinence, pressure injuries, or difficulty coping at home
    \item Concerns about capacity for decision-making or safe discharge from hospital
\end{itemize}

\section*{Differential Diagnosis}
A geriatrician considers many overlapping possibilities when assessing an older person's symptoms.

\begin{itemize}
    \item Delirium versus dementia versus depression in a confused older person
    \item Treatable causes of cognitive decline, including vitamin deficiencies, thyroid disease, and medicine effects
    \item Cardiac, neurological, and musculoskeletal causes of falls and dizziness
    \item Urinary tract infection and other acute illness presenting atypically in older people
    \item Adverse effects of medicines versus progression of underlying disease
    \item Social factors, including isolation, neglect, and carer strain, contributing to decline
\end{itemize}

\section*{When to Seek Medical Care}
Older people or their families should speak with their GP if there are concerns about memory, repeated falls, poor mobility, difficulty managing medicines, or decline in ability to cope at home. The GP can arrange referral to a geriatrician when specialist assessment is needed. Anyone can request a referral discussion with their doctor.

\textbf{Contact your local emergency services for:}
\begin{itemize}
    \item Sudden confusion, weakness on one side, facial drooping, or difficulty speaking, which may indicate stroke
    \item A fall with a suspected fracture or head injury
    \item Collapse, fainting, or loss of consciousness
    \item Severe breathlessness, chest pain, or other signs of serious acute illness
\end{itemize}

\section*{Diagnosis and Investigations}
The geriatrician's main diagnostic tool is the comprehensive geriatric assessment, supported by targeted tests.

\begin{itemize}
    \item \textbf{Comprehensive geriatric assessment:} a detailed evaluation of medical conditions, function, cognition, mood, continence, nutrition, and social supports
    \item \textbf{Cognitive testing:} standardised tools such as the Mini-Mental State Examination or Montreal Cognitive Assessment, plus specialist neuropsychological testing when needed
    \item \textbf{Mobility and balance assessment:} tests such as the timed up-and-go, gait observation, and falls risk screening
    \item \textbf{Blood tests and imaging:} to identify treatable causes of symptoms, including thyroid, vitamin B12, and brain imaging
    \item \textbf{Medicine review:} structured review of all prescribed and over-the-counter medicines
    \item \textbf{Functional assessment:} evaluation of the ability to perform activities of daily living
    \item \textbf{Home assessment:} an occupational therapist or community team may assess safety in the home environment
\end{itemize}

\section*{Management}
Geriatricians do not manage patients alone; they coordinate a multidisciplinary plan tailored to the individual's goals.

\begin{itemize}
    \item \textbf{Multidisciplinary team care:} physiotherapy, occupational therapy, nursing, pharmacy, and social work input
    \item \textbf{Medicine optimisation:} simplifying regimens, stopping unnecessary medicines, and avoiding prescribing cascades
    \item \textbf{Falls prevention:} strength and balance exercise, home safety modifications, and treatment of contributing conditions
    \item \textbf{Dementia care:} diagnosis, counselling, medication where appropriate, and support for carers
    \item \textbf{Rehabilitation:} hospital-based or community rehabilitation to restore function after illness or injury
    \item \textbf{Frailty management:} nutrition, exercise, and proactive management of chronic disease
    \item \textbf{Advance care planning:} helping people document their values and preferences for future care
    \item \textbf{Coordination of community services:} linking people with home care, respite, and aged care services
\end{itemize}

\section*{Complications}
Poorly managed ageing and multimorbidity are associated with several complications that geriatric care aims to prevent.

\begin{itemize}
    \item Falls, fractures, and head injuries
    \item Hospitalisation and hospital-acquired complications, including delirium and functional decline
    \item Adverse medicine events from polypharmacy
    \item Pressure injuries and incontinence
    \item Malnutrition and frailty
    \item Social isolation, depression, and carer burnout
    \item Premature admission to residential aged care when community supports fail
\end{itemize}

\section*{Prognosis and Outcomes}
Geriatric assessment and coordinated care improve outcomes for older people. Studies show that comprehensive geriatric assessment in hospital reduces functional decline, shortens stays, and reduces the need for residential care. Falls prevention programs reduce falls by around a quarter. The geriatrician's role is not to reverse ageing but to optimise health, independence, and quality of life, and to ensure that care matches the person's goals. With good geriatric care, many older people live well at home for longer despite complex health needs.

\section*{Prevention and Self-Care}
\begin{itemize}
    \item Maintain regular physical activity, including strength and balance exercise
    \item Eat well and maintain a healthy weight, with attention to protein and vitamin D
    \item Attend recommended health checks and screening for your age group
    \item Have medicines reviewed regularly by your GP and pharmacist
    \item Keep your home safe: remove trip hazards, improve lighting, and use handrails
    \item Stay socially connected and mentally active
    \item Plan ahead with advance care planning and discuss preferences with family and doctors
    \item Ask your GP for a referral to a geriatrician if you or a loved one has complex or deteriorating health
\end{itemize}

\section*{Sources and Further Reading}
The following reputable sources provide further information for healthcare students and patients seeking reliable, current advice about the role of geriatricians and care of older people.

\begin{itemize}
    \item Australian and New Zealand Society for Geriatric Medicine. \textit{About geriatric medicine.} https://www.anzsgm.org/
    \item Healthdirect Australia. \textit{Geriatrician: what they do and when to see one.} https://www.healthdirect.gov.au/
    \item Royal Australasian College of Physicians. \textit{Geriatric medicine: a career and role overview.} https://www.racp.edu.au/
    \item NHS. \textit{Geriatric medicine: what geriatricians do.} https://www.nhs.uk/
    \item American Geriatrics Society. \textit{Health in ageing foundation.} https://www.healthinaging.org/
    \item World Health Organization. \textit{Ageing and health.} https://www.who.int/
\end{itemize}
